\clearpage
\begin{trapbox}
\textbf{易错 1｜把“加 $a$ 变号”误当“加 $a$ 不变”}
\begin{center}\begin{tikzpicture}[font=\small,>=stealth,line join=round]
\draw[gray!40,rounded corners=3pt] (0,0) rectangle (14.2,8.3);\node at (7.1,7.700000000000001) {隔2变为相反高度，隔4才回到原高度};
\begin{scope}[shift={(1.7153846153846155,3.8842857142857143)},x=0.7692307692307693cm,y=1.3142857142857143cm]\draw[->,gray] (-0.8,0)--(7,0) node[right]{$x$};\draw[->,gray] (0,-1.7)--(0,1.8) node[above]{$y$};\draw[blue,thick,solid,domain=-0.4:6.6,samples=180] plot(\x,{sin(deg(pi*\x/2))});\fill[orange!85!black] (1,1) circle(2.2pt);\fill[orange!85!black] (3,-1) circle(2.2pt);\fill[orange!85!black] (5,1) circle(2.2pt);\foreach \v/\h in {1/1,3/-1,5/1}{\draw[gray,dashed] (\v,0)--(\v,\h);\node[below=5pt] at (\v,0){$\v$};}\end{scope}\node[anchor=west,align=left,text width=5.0cm] at (8.4,4.0) {$f=\sin\frac{\pi x}2,\ a=2$\\[9pt]$f(1)=1$\\[9pt]$f(3)=-1$\\[9pt]$f(5)=1$\\[9pt]正周期可取4};
\node[text width=12.6cm,align=center] at (7.1,.65) {蓝线是满足条件的一条示例曲线；证明只用任意输入的两次变号。};
\end{tikzpicture}\end{center}



错误：$f(x+a)=-f(x)$，所以 $a$ 是周期。\\
正确：加 $a$ 后函数值一般\textbf{不相等}，要再加一次，得到 $f(x+2a)=f(x)$。只有恒为零的特殊函数，$|a|$ 才也可能是周期。
\begin{center}\begin{tikzpicture}[font=\small,>=stealth,line join=round]
\draw[gray!40,rounded corners=3pt] (0,0) rectangle (14.2,8.3);\node at (7.1,7.700000000000001) {恒零函数：每一步既是反号，也同时保持原值};
\begin{scope}[shift={(2.317391304347826,3.8842857142857143)},x=0.8695652173913043cm,y=1.3142857142857143cm]\draw[->,gray] (-1.4,0)--(5.5,0) node[right]{$x$};\draw[->,gray] (0,-1.7)--(0,1.8) node[above]{$y$};\draw[blue,thick,solid,domain=-1:5,samples=180] plot(\x,{0});\fill[orange!85!black] (0,0) circle(2.2pt);\fill[orange!85!black] (2,0) circle(2.2pt);\fill[orange!85!black] (4,0) circle(2.2pt);\end{scope}\node[anchor=west,align=left,text width=5.0cm] at (8.4,4.0) {$f(x)\equiv0$\\[9pt]$-f(x)=f(x)=0$\\[9pt]任意正数都是周期\\[9pt]不存在最小正周期\\[9pt]此时 $|a|$ 也是周期};
\node[text width=12.6cm,align=center] at (7.1,.65) {非恒零函数移动|a|会在某处改变函数值；恒零函数是这里的例外。};
\end{tikzpicture}\end{center}



\medskip
\end{trapbox}\clearpage\begin{trapbox}\textbf{易错 2｜把“一个周期”写成“最小正周期”}
\begin{center}\begin{tikzpicture}[font=\small,>=stealth,line join=round]
\draw[gray!40,rounded corners=3pt] (0,0) rectangle (14.2,8.3);\node at (7.1,7.700000000000001) {例4：变号步长是π，重复却早在三分之二π出现};
\begin{scope}[shift={(1.4287671232876713,3.8842857142857143)},x=0.821917808219178cm,y=1.3142857142857143cm]\draw[->,gray] (-0.4,0)--(6.9,0) node[right]{$x$};\draw[->,gray] (0,-1.7)--(0,1.8) node[above]{$y$};\draw[blue,thick,solid,domain=-0.1:6.5,samples=180] plot(\x,{sin(deg(3*\x))});\fill[orange!85!black] (0.5235987755982988,1) circle(2.2pt);\fill[orange!85!black] (2.6179938779914944,1) circle(2.2pt);\fill[orange!85!black] (4.71238898038469,1) circle(2.2pt);\foreach \v/\lab in {2.094395/\frac{2\pi}3,4.18879/\frac{4\pi}3,6.283185/2\pi}{\node[below=5pt] at (\v,-1.25){$\lab$};}\end{scope}\node[anchor=west,align=left,text width=5.0cm] at (8.4,4.0) {$f=\sin3x,\ a=\pi$\\[9pt]$f(x+\pi)=-f(x)$\\[9pt]$2\pi$ 是一个周期\\[9pt]最小正周期 $\frac{2\pi}{3}$};
\node[text width=12.6cm,align=center] at (7.1,.65) {三个波形在长度2π内重复；两次变号只保证得到一个周期。};
\end{tikzpicture}\end{center}



由题设只能保证 $2|a|$ \textbf{是一个正周期}。例如 $f(x)=\sin3x$ 与 $a=\pi$ 相配，最小正周期是 $2\pi/3$，而不是 $2\pi$。

\medskip
\end{trapbox}\clearpage\begin{trapbox}\textbf{易错 3｜从变号关系猜出一个不存在的对称中心（原稿重要修正）}
\begin{center}\begin{tikzpicture}[font=\small,>=stealth,line join=round]
\draw[gray!40,rounded corners=3pt] (0,0) rectangle (14.2,8.5);\node at (7.1,7.9) {同样加π变号，额外的奇偶性决定不同对称};
\node at (2.4,6.9) {$f=\sin x$：奇函数};\begin{scope}[shift={(1.2666666666666666,4.341935483870968)},x=0.7083333333333334cm,y=1.1612903225806452cm]\draw[->,gray] (-0.8,0)--(4,0) node[right]{$x$};\draw[->,gray] (0,-1.5)--(0,1.6) node[above]{$y$};\draw[blue,thick,solid,domain=-0.5:3.8,samples=160] plot(\x,{sin(deg(\x))});\fill[orange!85!black] (0.5235987755982988,0.5) circle(2pt);\fill[orange!85!black] (2.6179938779914944,0.5) circle(2pt);\draw[orange!85!black,dashed] (1.570796,-1.2)--(1.570796,1.3);\end{scope}\node[align=center,text width=4cm] at (2.4,1.7) {关于直线 $x=\frac\pi2$ 对称};\node at (7.050000000000001,6.9) {$g=\cos x$：偶函数};\begin{scope}[shift={(5.916666666666667,4.341935483870968)},x=0.7083333333333334cm,y=1.1612903225806452cm]\draw[->,gray] (-0.8,0)--(4,0) node[right]{$x$};\draw[->,gray] (0,-1.5)--(0,1.6) node[above]{$y$};\draw[violet,thick,solid,domain=-0.5:3.8,samples=160] plot(\x,{cos(deg(\x))});\fill[orange!85!black] (0.5235987755982988,0.8660254037844386) circle(2pt);\fill[orange!85!black] (2.6179938779914944,-0.8660254037844386) circle(2pt);\fill[orange!85!black] (1.570796,0) circle(2.2pt);\draw[orange!85!black,dashed] (.523599,.866025)--(2.617994,-.866025);\end{scope}\node[align=center,text width=4cm] at (7.050000000000001,1.7) {关于点 $(\frac\pi2,0)$\\中心对称};\node at (11.7,6.9) {$f=\sin x$ 的中心检验};\begin{scope}[shift={(10.566666666666666,4.341935483870968)},x=0.7083333333333334cm,y=1.1612903225806452cm]\draw[->,gray] (-0.8,0)--(4,0) node[right]{$x$};\draw[->,gray] (0,-1.5)--(0,1.6) node[above]{$y$};\draw[blue,thick,solid,domain=-0.5:3.8,samples=160] plot(\x,{sin(deg(\x))});\fill[orange!85!black] (1.5707963267948966,1) circle(2pt);\end{scope}\node[align=center,text width=4cm] at (11.7,1.7) {$f(\frac\pi2)=1\ne0$\\$(\frac\pi2,0)$ 不是中心};
\node[text width=12.6cm,align=center] at (7.1,.65) {三图同尺度；只有变号条件时，不能直接指定对称轴或对称中心。};
\end{tikzpicture}\end{center}



错误：$f(x+a)=-f(x)$，所以图像一定关于 $\bigl(\frac a2,0\bigr)$ 中心对称。\\
\textbf{反例：}$f(x)=\sin x$、$a=\pi$，虽然 $\sin(x+\pi)=-\sin x$，但 $f(\pi/2)=1\ne0$。若图像真关于 $(\pi/2,0)$ 中心对称，取正中央横坐标时应有 $f(\pi/2)=0$，因此矛盾。\\
\textbf{怎样才能说对称？}额外知道\textbf{偶函数}，才能推出上述对称中心；额外知道\textbf{奇函数}，则推出 $x=a/2$ 是对称轴。

\medskip
\end{trapbox}\clearpage\begin{trapbox}\textbf{易错 4｜只对某几个数字成立，就当成函数的周期}
\begin{center}\begin{tikzpicture}[font=\small,>=stealth,line join=round]
\draw[gray!40,rounded corners=3pt] (0,0) rectangle (14.2,8.3);\node at (7.1,7.700000000000001) {单独一对反号，无法推出整幅图像的周期};
\begin{scope}[shift={(2.363157894736842,3.5666666666666664)},x=1.0526315789473684cm,y=0.7666666666666666cm]\draw[->,gray] (-1.2,0)--(4.5,0) node[right]{$x$};\draw[->,gray] (0,-2.5)--(0,3.5) node[above]{$y$};\draw[blue,thick,solid,domain=-0.8:4,samples=180] plot(\x,{\x-1});\fill[orange!85!black] (0,-1) circle(2.2pt);\fill[orange!85!black] (2,1) circle(2.2pt);\fill[orange!85!black] (1,0) circle(2.2pt);\fill[orange!85!black] (3,2) circle(2.2pt);\end{scope}\node[anchor=west,align=left,text width=5.0cm] at (8.4,4.0) {$f=x-1,\ a=2$\\[9pt]$f(0)=-1,\ f(2)=1$\\[9pt]这一对互为相反数\\[9pt]但 $f(1)=0,\ f(3)=2$\\[9pt]全称条件不成立};
\node[text width=12.6cm,align=center] at (7.1,.65) {周期关系必须对所有输入成立；不能从两点推断整条直线会重复。};
\end{tikzpicture}\end{center}



周期要求对应的等式对\textbf{所有允许的 $x$} 都成立。只知道 $f(1+a)=-f(1)$，不能推断整个函数都是如此。若定义域不是全体实数，还必须保证 $x$ 和移动后的输入都处在定义域中；否则不能随意继续“加一次、再加一次”。

\medskip
\end{trapbox}\clearpage\begin{trapbox}\textbf{易错 5｜负数 $a$ 时写“正周期 $T=2a$”}
\begin{center}\begin{tikzpicture}[font=\small,>=stealth,line join=round]
\draw[gray!40,rounded corners=3pt] (0,0) rectangle (14.2,8.3);\node at (7.1,7.700000000000001) {向左走3变号，向左走6还原；正周期取6};
\begin{scope}[shift={(4.1,3.8842857142857143)},x=0.42857142857142855cm,y=1.3142857142857143cm]\draw[->,gray] (-7,0)--(7,0) node[right]{$x$};\draw[->,gray] (0,-1.7)--(0,1.8) node[above]{$y$};\draw[blue,thick,solid,domain=-6.5:6.5,samples=180] plot(\x,{cos(deg(pi*\x/3))});\fill[orange!85!black] (-6,1) circle(2.2pt);\fill[orange!85!black] (-3,-1) circle(2.2pt);\fill[orange!85!black] (0,1) circle(2.2pt);\fill[orange!85!black] (6,1) circle(2.2pt);\foreach \v in {-6,-3,0,6}{\ifnum\v=0\node[below left=5pt] at (\v,0){$\v$};\else\node[below=5pt] at (\v,0){$\v$};\fi}\end{scope}\node[anchor=west,align=left,text width=5.0cm] at (8.4,4.0) {图取 $a=-3$\\[9pt]$f=\cos\frac{\pi x}3$\\[9pt]$f(-3)=-f(0)$\\[9pt]$f(-6)=f(0)=f(6)$\\[9pt]$T=6=2|a|$};
\node[text width=12.6cm,align=center] at (7.1,.65) {带符号的位移可以是负数；正周期表示正的重复距离。};
\end{tikzpicture}\end{center}



若 $a=-3$，推导出来的是 $f(x-6)=f(x)$，可以把两边的 $x$ 换成 $x+6$，得 $f(x+6)=f(x)$。\textbf{正周期应写 $6=2|a|$。}

\medskip
\end{trapbox}\clearpage\begin{trapbox}\textbf{易错 6｜例题在两个端点给出了冲突的数值}
\begin{center}\begin{tikzpicture}[font=\small,>=stealth,line join=round]
\draw[gray!40,rounded corners=3pt] (0,0) rectangle (14.2,8.3);\node at (7.1,7.700000000000001) {例2：每段长3，下一段反号；接头不能同时取两值};
\begin{scope}[shift={(1.3891566265060242,3.828947368421052)},x=0.36144578313253006cm,y=0.4842105263157894cm]\draw[->,gray] (-0.8,0)--(15.8,0) node[right]{$x$};\draw[->,gray] (0,-4.5)--(0,5) node[above]{$y$};\draw[blue,thick,solid,domain=0:3,samples=180] plot(\x,{1*(\x-0)});\draw[blue,thick,solid,domain=3:6,samples=180] plot(\x,{-1*(\x-3)});\draw[blue,thick,solid,domain=6:9,samples=180] plot(\x,{1*(\x-6)});\draw[blue,thick,solid,domain=9:12,samples=180] plot(\x,{-1*(\x-9)});\draw[blue,thick,solid,domain=12:15,samples=180] plot(\x,{1*(\x-12)});\fill[orange!85!black] (0,0) circle(2.2pt);\fill[orange!85!black] (3,0) circle(2.2pt);\fill[orange!85!black] (6,0) circle(2.2pt);\fill[orange!85!black] (9,0) circle(2.2pt);\fill[orange!85!black] (12,0) circle(2.2pt);\fill[orange!85!black] (1,1) circle(2.2pt);\fill[orange!85!black] (10,-1) circle(2.2pt);\fill[orange!85!black] (13,1) circle(2.2pt);\draw[blue,fill=white] (3,3) circle(2.8pt);\draw[blue,fill=white] (6,-3) circle(2.8pt);\draw[blue,fill=white] (9,3) circle(2.8pt);\draw[blue,fill=white] (12,-3) circle(2.8pt);\draw[blue,fill=white] (15,3) circle(2.8pt);\fill[orange!85!black] (15,0) circle(2.2pt);\end{scope}\node[anchor=west,align=left,text width=5.0cm] at (8.4,4.0) {$[0,3)$ 上 $f=x$\\[9pt]$f(x+3)=-f(x)$\\[9pt]$f(1)=1,\ f(10)=-1$\\[9pt]$f(13)=1$\\[9pt]实点取到，空圈不取};
\node[text width=12.6cm,align=center] at (7.1,.65) {每段左端取0，右端的正负3不取；否则与变号条件发生冲突。};
\end{tikzpicture}\end{center}



若条件为 $f(x+3)=-f(x)$，就必须有 $f(3)=-f(0)$。因此“在 $[0,3]$ 上 $f(x)=x$”不成立：前者给 $f(0)=0$、$f(3)=3$，后者却要求 $f(3)=0$。可将已知区间改成 $[0,3)$，再由关系补齐其他部分。

\medskip
\end{trapbox}\clearpage\begin{trapbox}\textbf{易错 7｜连续才有“过横轴”}
\begin{center}\begin{tikzpicture}[font=\small,>=stealth,line join=round]
\draw[gray!40,rounded corners=3pt] (0,0) rectangle (14.2,8.3);\node at (7.1,7.700000000000001) {跳跃函数也能每次反号，却始终不取0};
\begin{scope}[shift={(2.7595744680851064,3.8842857142857143)},x=1.276595744680851cm,y=1.3142857142857143cm]\draw[->,gray] (-1.3,0)--(3.4,0) node[right]{$x$};\draw[->,gray] (0,-1.7)--(0,1.8) node[above]{$y$};\draw[blue,thick,solid,domain=-1:0,samples=180] plot(\x,{-1});\draw[blue,thick,solid,domain=0:1,samples=180] plot(\x,{1});\draw[blue,thick,solid,domain=1:2,samples=180] plot(\x,{-1});\draw[blue,thick,solid,domain=2:3,samples=180] plot(\x,{1});\fill[orange!85!black] (-1,-1) circle(2.2pt);\fill[orange!85!black] (0,1) circle(2.2pt);\fill[orange!85!black] (1,-1) circle(2.2pt);\fill[orange!85!black] (2,1) circle(2.2pt);\draw[blue,fill=white] (0,-1) circle(2.8pt);\draw[blue,fill=white] (1,1) circle(2.8pt);\draw[blue,fill=white] (2,-1) circle(2.8pt);\draw[blue,fill=white] (3,1) circle(2.8pt);\end{scope}\node[anchor=west,align=left,text width=5.0cm] at (8.4,4.0) {$f(x)=(-1)^{\lfloor x\rfloor}$\\[9pt]$\lfloor x\rfloor$ 表示不超过\\$x$ 的最大整数\\[9pt]$f(x+1)=-f(x)$\\[9pt]函数值只有 $1,-1$};
\node[text width=12.6cm,align=center] at (7.1,.65) {实点取到，空圈不取；这条函数在整数处跳跃，不连续，也没有零点。};
\end{tikzpicture}\end{center}



若 $a>0$ 且 $f$ 在 $[x,x+a]$ 上连续，两端数值相反，才可以说曲线至少会碰到一次 $x$ 轴。若图像可以突然跳开，即使两端一正一负，也未必在中间取到 $0$。不能偷偷漏掉“连续”。
\end{trapbox}
